

\noindent The equation we are solving is the following:
\begin{equation}
 i\frac{\partial}{\partial t}\psi(x,t) = 
	\left(-\frac{1}{2}\frac{\partial^2}{\partial x^2} 
           + \frac{1}{2}\omega^2x^2\right)\psi(x,t),
	  \hspace{9mm} \omega=1 \label{eqn:tdse}
\end{equation}

\noindent with the fundamental as the inital condition, 
unless stated otherwise.

\section{Analytical solution and convergence control}
\subsection{Time Independent Schr\"odinger Equation}

\noindent We first solve the TISE given by 
\begin{equation}
	-\frac{1}{2}\frac{\partial^2}{\partial x^2}\varphi(x) 
	  + \frac{1}{2}\omega^2x^2\varphi(x) = E\varphi(x), \label{eqn:tise}
\end{equation}

\noindent with the following boundary conditions
\begin{equation}
		\varphi(x) \xrightarrow{\tiny x \rightarrow \pm\infty} 0.
\end{equation}

\noindent The static solutions are given by the following equation,
\begin{equation}
 \varphi_n(x) = 
	\left(\frac{\omega}{\pi}\right)^{1/4}
	 \frac{1}{\sqrt{2^n n!}}H_n(\sqrt{\omega}x)e^{-\omega x^2/2}, 
	 \label{eqn:eigenfun}
\end{equation}

\noindent whose the eigenvalues reads 

\begin{equation}
	E_n = \omega\left(n + \frac{1}{2} \right). \label{eqn:eigenval}
\end{equation}

\noindent In the following, we set $\omega=1$. The fundamental $(n=0)$ will be
\begin{equation}
	\varphi_0(x) = \left(\pi\right)^{-1/4} e^{-x^2/2} \label{eqn:init}
\end{equation}


\subsection{Time Dependent Schr\"odinger Equation}

We want to solve (\ref{eqn:tdse}), 
with the initial state given by (\ref{eqn:init}). 
Since 
${H}= -\frac{1}{2}\frac{\partial^2}{\partial x^2} + \frac{1}{2}x^2$ 
does not depend on time, the solution is simply given by

\begin{align}
	\psi(x,t) &  = \varphi_0(x)e^{-iHt}  \\
	 & = \left(\pi\right)^{-1/4} e^{-\frac{1}{2}(x^2+it)}  
                   \label{eqn:tdse_sol}
\end{align}

\noindent Additionally, we can verify that the following general expression
is a solution to (\ref{eqn:tdse})
\begin{\eq}
	\psi(x,t) = \left(\pi\right)^{-1/4} 
           \exp \left[-\frac{1}{2}\left(x^2+
                \frac{a^2}{2}\left(1+e^{-2it}\right) 
                + it -2axe^{-it}\right) \right],
                   \label{eqn:tdse_gensol_1}
\end{\eq}

\noindent where $a$ can be any real parameter with the dimension of length.
Let's note that for $a=0$, we have the basic solution (\ref{eqn:tdse_sol}).
We may rewrite this last general expression as the following

\begin{\eq}
	\psi(x,t) = \left(\pi\right)^{-1/4}
         \exp\left\lbrace -\frac{1}{2}
\left[ \left(x-a\cos t\right)^2 
               +i\left( t +2ax\sin t 
               - \frac{a^2}{2}\sin 2t\right)
\right] 
\right\rbrace.
                  \label{eqn:tdse_gensol_2}
\end{\eq}

\newpage

\subsection{TDSE for $V(x)=\omega^2(x+1)^2/2$}

We want to solve the Schr\"odinger equation 
\begin{equation}
 i\frac{\partial}{\partial t}\psi(x,t) = 
	\left(-\frac{1}{2}\frac{\partial^2}{\partial x^2} 
          + \frac{1}{2}\omega^2(x+1)^2\right)\psi(x,t),
    \label{eqn:tdse_2}
\end{equation}

\noindent and the following initial condition

\begin{equation}
	\Psi(x,0) = \left(\pi\right)^{-1/4} e^{-x^2/2} \label{eqn:init2}.
\end{equation}

\noindent We can infer the solution from (\ref{eqn:tdse}), 
by making the substitution
\begin{equation}
   x \rightarrow y+1   \label{eqn:change_var}
\end{equation}

\noindent We then have the following
\begin{align}
    dx = dy \hspace{5mm} \Longrightarrow \hspace{5mm} 
	   \frac{\partial}{\partial x}  = \frac{\partial}{\partial y}, \\
	x \rightarrow \infty \hspace{5mm} \Longrightarrow \hspace{5mm} 
	 y \rightarrow \infty,
\end{align}

\noindent and may rewrite the Schr\"odinger equation (\ref{eqn:tdse})

\begin{equation}
 i\frac{\partial}{\partial t}\Psi(y,t) = 
	\left(-\frac{1}{2}\frac{\partial^2}{\partial y^2} + \frac{1}{2}\omega^2(y+1)^2\right)\Psi(y,t),
	  \hspace{9mm} \omega=1 \label{eq:tdse2}.
\end{equation}

\noindent The static solutions are still given by (\ref{eqn:eigenfun}), 
that we may rewrite in terms of $y$ and with $\omega=1$, as

\begin{equation}
 \varphi_n(y) = 
	\left(\frac{1}{\pi}\right)^{1/4}
	 \frac{1}{\sqrt{2^n n!}}H_n(y+1)e^{-(y+1)^2/2}.
	 \label{eq:static2}
\end{equation}

\noindent The solution to the time dependent problem is then given by 
(\ref{eqn:tdse_gensol_2}), provided the substitution (\ref{eqn:change_var})

\begin{\eq}
	\psi(x,t) = \left(\pi\right)^{-1/4}
         \exp\left\lbrace -\frac{1}{2}
\left[ \left(x+1-a\cos t\right)^2 
               +i\left( t +2a(x+1)\sin t 
               - \frac{a^2}{2}\sin 2t\right)
\right] 
\right\rbrace.
                  \label{eqn:tdse_gensol_3}
\end{\eq}

This is the most general solution to 
the Schr\"odinger equation (\ref{eqn:tdse_2}). The initial condition
(\ref{eqn:init2}) is satisfied by setting $a=1$.


\vspace{2cm}

%\noindent In practice, 
%an easy implementation of (\ref{eq:tdse_sol2}) in the current code
%can be done by making the substitution $x \rightarrow x - 1$. Doing so, we 
%would use the following solution
%\begin{align}
%	\Psi(x,t) &  = \varphi_0(x-1)e^{-iHt},  \\
%	 & = \left(\pi\right)^{-1/4} e^{-\frac{1}{2}\left((x-1)^2+it\right)}. 
%\end{align}

%\newpage
%
%\noindent This is the implementation in the script
%
%\begin{verbatim}
%DT = 0.01d0
%DT2 = DT/2.d0
%NTIME = 10000
%DO ITIME=1,NTIME                 ! This loop should start at 1, not 0
%   TIME = ITIME*DT
%
%C
%C
%C
%
%   DO  j= 1,20
%      IF (j.EQ.1) THEN 
%         HN(j) = 1.d0
%      ELSE If (j.EQ.2) THEN
%         HN(j) = 2*XX(II)
%      ELSE
%         HN(j) = 2*XX(II)*HN(j-1) * 2*(j-1)*HN(j-2)
%   ENDDO
%
%   n = 0                       ! Parameter in \varphi_n
%   j = n + 1
%   DO II = 1, NX
%      CE0 = DCMPLEX(0.d0,.5d0)
%      C4(II) = CDEXP(-CE0*DT)
%      C4(II) = HN(j)*DEXP(-.5*XX(II)**2)*CDEXP(-CE0*DT)*C4(II)
%      C4(II) = (1/DSQRT(2**n fact(n))*C4(II)
%      C4(II) = (4.d0*DATAN(1.d0))**(.25)*C4(II)
%
%      WRITE(99,) TIME, X(II), C3(II)/WX(II)/dsqrt(hscale), C4(II)
%      WRITE(77,) TIME, X(II), CDABS(C4(II))
%
%C  ----- Evaluation of the error
%      CERR1 = C3(II)/WX(II)/dsqrt(hscale) - C4(II)
%      CERR2 = CERR1/CDABS(C4(II))
%      CERR1 = CDABS(CERR1)                ! Absolute error
%      CERR2 = CDABS(CERR2)                ! Relative error
%      WRITE(88,*) TIME, X(II), CERR1, CERR2 
%
%   ENDDO
%ENDDO
%\end{verbatim}
%
